\documentclass[10pt,a4paper]{amsart}
\usepackage[margin=19mm]{geometry}
\usepackage{amsmath,amssymb,mathtools}
\usepackage[scr=rsfs,cal=euler]{mathalfa}
\usepackage{xfrac}
\usepackage[unicode,hidelinks,bookmarks=false]{hyperref}
\newcommand{\ie}{i.e.\ }
\newcommand{\cf}{cf.\ }
\newcommand{\resp}{respectively\ }
\newcommand{\Subsection}{\subsection}
\providecommand{\nobreakdash}{\nobreak\hbox{-}}
\setlength{\emergencystretch}{2em}
\multlinegap0pt
\usepackage{tikz}
\usepackage{bm}
\usepackage{amssymb}
\usepackage{enumerate}
\usepackage[normalem]{ulem}
\newtheorem{lemma}{Lemma}
\newtheorem{claim}{Claim}
\newtheorem{obser}{Observation}
\def\andrey#1{{\sc Andrey: }{\marrow\sf \textcolor{blue}{#1}}}
\newcommand{\m}{\mathcal}
\def\marrow{{\boldmath {\marginpar[\hfill$\rightarrow \rightarrow$]{$\leftarrow \leftarrow$}}}}
\title{Structure and properties of large cross-intersecting families}
\author{Yang Huang, Andrey Kupavskii}
\date{Source: arXiv:2606.20085v1 (CC BY 4.0)}
\begin{document}
\maketitle

\noindent\emph{Source and licence.} Structure and properties of large cross-intersecting families. Version arXiv:2606.20085v1 (CC BY 4.0), distributed by arXiv under the Creative Commons Attribution 4.0 International licence, http://creativecommons.org/licenses/by/4.0/, as recorded in the arXiv licence metadata for this exact version and shown on the abstract page https://arxiv.org/abs/2606.20085v1. Full licence notice: \texttt{licenses/arxiv-2606.20085-CC-BY-4.0.txt}.\par\smallskip

\noindent\textit{Editorial scope.} Counted unit: the article's claim 11-19-1 (source0.tex lines 2498--2500). The article's complete original proof of it is delivered with it (source0.tex lines 2503--2769, one \texttt{proof} environment of 267 lines). No mathematics is written, repaired or re-derived: the statement and the proof are the article's own lines, in the article's own notation, and no retained line is substituted or re-flowed.  Retained uncounted as the source-local context the proof needs: lines 541--544; lines 1807--1810; lines 1839--1841; lines 2006--2015; lines 2130--2137; lines 2332--2334.  Nothing was omitted from any retained span, so the package carries no omission record.  No retained line contains a citation, so the wrapper carries no bibliography and no bibliography entry.  No retained line contains a figure environment, an image-inclusion command or drawing code, so no image is delivered and the package declares no asset dependency.  Editorial formatting only, in addition: an AMS \texttt{amsart} preamble that restates the article's own declarations and macros used by the retained lines, namely the environments \texttt{lemma}, \texttt{claim}, \texttt{obser} on separate counters, together with the standard packages those lines need.  The retained environments print in this document as Lemma 1, Claim 1, Obser 1 in order of appearance, whereas the article prints the same items with the numbering of its own full document.  The complete native source of the article as distributed in the arXiv e-print for this exact version is bundled unchanged as \texttt{sources/203192/source0.tex}.
\begin{lemma}\label{SUV}
Let $n\ge a+b$, \(\mathcal{A}\subset {[n]\choose a}\) and \(\mathcal{B}\subset {[n]\choose b}\) be cross-intersecting.
Let $(U,V)$ satisfy {\bf P} for $(\m A, \m B)$. Then \(S_{U,V}(\mathcal{A})\) and \(S_{U,V}(\mathcal{B})\) are also cross-intersecting.
\end{lemma}

\begin{lemma}\label{SUV'}
Let $n\ge a+b$, \(\mathcal{A}\subset {[n]\choose a}\) and \(\mathcal{B}\subset {[n]\choose b}\) be cross-intersecting. 
Let $(U,V)$ satisfy {\bf Q} for $(\m A, \m B)$. Then \(S_{U,V}^{Q}(\mathcal{A})\) and \(S_{U,V}^{Q}(\mathcal{B})\) are also cross-intersecting.
\end{lemma}

\begin{equation}\label{ec}
    \mathcal{E}^k_C:=\{C\cup T: T\in \tbinom{[\max C+1, n]}{ k-|C|}\}.
\end{equation}

\begin{claim}\label{11-3-4}
Let $\mathcal{C}\subset{[n]\choose a}, \mathcal{D}\subset {[n]\choose b}$ be cross-intersecting families satisfying  $|\mathcal{C}|= |\mathcal{A}|$, $|\mathcal{D}|= |\mathcal{B}|$,
$|\m C(1)|\ge |\m A(1)|$, $|\m D(1)|\ge |\m B(1)|$,
%$\Delta(\mathcal{C})=|\mathcal{C}(1)|\ge {n-2\choose a-2}+\gamma(\mathcal{C})$, $\Delta(\mathcal{D})=|\mathcal{D}(1)|\ge {n-2\choose b-2}+\gamma(\mathcal{D})$
 and $(\mathcal{C}, \mathcal{D})\precneqq (\mathcal{A},\mathcal{B})$. 
 Then $\mathcal{M}\not\subset \mathcal{C}(\bar1)$ or $\mathcal{T}\not\subset \mathcal{D}(\bar1)$. 
 In particular, 
for any $S^Q_{U,V}$-shift  of $(\m A, \m B)$, 
we have $\mathcal{M}\not\subset S^Q_{U,V}(\mathcal{A})(\bar1)$ or $\mathcal{T}\not\subset S^Q_{U,V}(\mathcal{B})(\bar1)$.  
\end{claim}

\begin{claim}\label{abcd}
\begin{itemize}
\item[(i)] For each $A\in \mathcal{A}(\bar{1})$, we have $[2, m']\subset A$, and there exists $A'\in \mathcal{A}(\bar{1})$ such that $m'+1\not\in A'$.
\item[(ii)] For each $i \in [2, m']$, $\{1, i\}$ is full in $\mathcal{B}$.
\item[(iii)] $\{1, m'+1\}$ is non-full in $\mathcal{B}$.
\item[(iv)] For each $i\ge m'+2$, $\{1, i\}$ is empty in $\mathcal{B}$. 
\end{itemize}
\end{claim}

\begin{claim}\label{11-13-1}
Every set of $\mathcal{B}(\bar{1})\setminus \mathcal{T}$ contains $[2, t+1+\ell]$. 
\end{claim}

\begin{claim}\label{11-19-1}
Every set of $\mathcal{B}(\bar{1})\setminus \mathcal{T}$ contains $[2, m'+2]$.
\end{claim}

\begin{proof} 
Since Claim \ref{11-13-1}, every set in $\m B(\bar1) \setminus \m T$ contains $[2,t+1+\ell]$. 
If $t+\ell\ge m'+1$, then we are done. Assume 
\[
t+\ell\le m'
\]

\begin{obser}\label{obser3}
Let $B \in \m B(\bar1) \setminus \m T$, and let $U,V\subset [n]$ with $U\cap V=\emptyset$, $1\in U$ and $|U|=|V|$. If  $S_{U,V}(B)\not\in \m B$, then $[2,t+\ell]\subset V$ and $|U|=|V|\ge \ell+2$. 
\end{obser}

\begin{proof}
Since $B \in \m B(\bar1) \setminus \m T$, Claim \ref{11-13-1} gives $[2,t+1+\ell]\subset B$. 
Since $1\in U$ and $U\cap V=\emptyset$, $U\prec V$. 
%As $S_{U,V}(B)\not\in \m B$, $S_{U,V}(B)\precneqq B$ 
By Claim \ref{abcd}(ii), $S_{U,V}(B)\cap [2,m']=\emptyset$.
If $t+1+\ell\le m'$, then $[2,t+1+\ell]\subset V$. 
If $t+1+\ell\ge m'+1$, then since $t+\ell \le m'$, we have $m'=t+\ell$. Thus $[2,m'+1]=[2,t+1+\ell]\subset B$.
As $S_{U,V}(B)\cap [2,m']=\emptyset$, $[2, t+\ell]\subset V$. Since $t\ge 3$ and $|U|=|V|$, $|V|=|U|\ge \ell+2$.
\end{proof}

We will find a pair $(U,V)$ such that $(S_{U,V}(\m A), S_{U,V}(\m B\setminus \m T)\cup \m T)\precneqq (\m A, \m B)$, $|S_{U,V}(\m B\setminus \m T)\cup \m T|=|\m B|$,
$S_{U,V}(\m A)$ and $S_{U,V}(\m B\setminus \m T)\cup \m T$ are cross-intersecting, and $\m M $, $\m T$ are stable. This contradicts Claim \ref{11-3-4}, thereby completing the proof of Claim \ref{11-19-1}. 

We choose $U$ and $V$ according to the following cases.  
We shall later verify that this $S_{U,V}$-shift preserves $\m M, \m T$ and the cross-intersecting property.

For a family $\m F\subset {[n]\choose k}$ and a set $C\subset [n]$ with $|C|\le k$, we define
\begin{equation}\label{fc}
    \mathcal{F}^k_C:=\{C\cup R: R\in \tbinom{[n]\setminus C}{ k-|C|}\}.
\end{equation}
Then $\mathcal{E}^k_C\subset \mathcal{F}^k_C$ (the family $\mathcal{E}^k_C$ is defined in (\ref{ec})). Moreover, if $C$ is non-full in $\m F$, then $\mathcal{E}^k_C\setminus \m F \ne \emptyset$, and hence
$\mathcal{F}^k_C\setminus \m F \ne \emptyset$. %(\andrey{But the converse is not immediately true.}) 

\begin{itemize}
\item[Case 1.] There exist $B\in \m B(\bar{1})\setminus \m T$ and $I_B$ such that $m'+2\not\in B $ and  $\{1,m'+2\}\cup I_B$ is non-full in $\m B$. 

Fix $B_1$ that satisfies Case 1. 
%Then $\{1,m'+2\}\cup I_{B_1}$ is non-full in $\m B$. 
Let $R_1:=\{1,m'+2\}\cup I_{B_1}$.
Define 
\begin{equation*}
    \m F_1:=\{F\in \m B(\bar1): R_1\cap F=\emptyset\}, \,\,\,
    \m G_1:=\m F_{R_1}^b\setminus \m B.
\end{equation*}
Since $B_1\in \m F_1$, $\m F_1\ne \emptyset$.
Since $R_1$ is non-full in $\m B$, $\m G_1\ne \emptyset$. 
For any $F\in \m F_1$ and $G\in \m G_1$, there exists a pair $(U,V)$ such that $S_{U,V}(F)=G$. Among all such pairs $(U,V)$, choose one with $V$ inclusion-minimal. 


\item[Case 2.] 
There exist $B\in \m B(\bar{1})\setminus \m T$ and $I_B$ such that $m'+2\not\in B $ and $\{1,m'+2\}\cup I_B$ is full in $\m B$. 

Fix $B_2$ that satisfies Case 2.
%Then $\{1,m'+2\}\cup I_{B_2}$ is full in $\m B$. 
Define $I'_{B_2}$ as follows: 
if $\{1,m'+2\}\cup (I_{B_2}\setminus [2,m'])$ is non-full, then 
let $I'_{B_2}=I_{B_2}\setminus [2,m']$; otherwise, let $I'_{B_2}=I_{B_2}\setminus [2,m'+1]$. 
Note that $\{1,m'+2\}$ is empty in $\m B$. Thereby $\{1,m'+2\}\cup I'_{B_2}$ is non-full in $\m B$. Moreover, if $I'_{B_2}=I_{B_2}\setminus [2,m'+1]$, then $\{1,m'+1,m'+2\}\cup I'_{B_2}$ is full. 
Let $R_2:=\{1, m'+2\}\cup I'_{B_2}$.
Define 
\begin{equation*}
    \m F_2:=\{F\in \m B(\bar1): (\{m'+2\}\cup I_{B_2})\cap F=\emptyset\}, \,\,\,
    \m G_2:=\m F_{R_2}^b\setminus \m B.
\end{equation*}
Since $B_2\in \m F_2$, $\m F_2\ne \emptyset$. Since $R_2$ is non-full in $\m B$, $\m G_2\ne \emptyset$. 
For any $F\in \m F_2$ and $G\in \m G_2$, there exists a pair $(U,V)$ such that $S_{U,V}(F)=G$. Among all such pairs $(U,V)$, choose one with $V$ inclusion-minimal. 



\item[Case 3.] Every set in $\m B(\bar{1})\setminus \m T$ contains $m'+2$, and there exist $B\in \m B(\bar{1})\setminus \m T$ and $I_B$ such that $\{1,m'+2\}\cup I_B$ is non-full in $\m B$.

Fix $B_3$ that satisfies Case 3. 
%Then $\{1,m'+2\}\cup I_{B_3}$ is non-full in $\m B$.
Let $R_3:=\{1\}\cup I_{B_3}$.
Define
\begin{equation*}
    \m F_3:=\{F\in \m B(\bar1): R_3\cap F=\emptyset\}, \,\,\,
    \m G_3:=\m F_{R_3\cup \{m'+2\}}^b\setminus \m B.
\end{equation*}
Since $B_3\in \m F_3$, $\m F_3\ne \emptyset$.
Since $R_3\cup \{m'+2\}$ is non-full in $\m B$, $\m G_3\ne \emptyset$. 
For any $F\in \m F_3$ and $G\in \m G_3$, there exists a pair $(U,V)$ such that $S_{U,V}(F)=G$. Among all such pairs $(U,V)$, choose one with $V$ inclusion-minimal. 

\item[Case 4.] Every set in $\m B(\bar{1})\setminus \m T$ contains $m'+2$, and there exists $B\in \m B(\bar{1})\setminus \m T$
such that $\{1,m'+2\}\cup I_B$ is full in $\m B$. 
Furthermore, assume that $\{1\}\cup I_{B}$ is non-full in $\m B$.

Fix $B_4$ that satisfies Case 4. 
%Then $\{1,m'+2\}\cup I_{B_4}$ is full in $\m B$, and $\{1\}\cup I_{B_4}$ is non-full in $\m B$. 
%Note that $\{1,m'+2\}$ is empty in $\m B$ (see Cliam \ref{abcd}(iv)). Then since $\{1,m'+2\}\cup I_{B_4}$ is full in $\m B$, $I_{B_4}\cap [2,m'+1]\ne \emptyset$; and since $\{1\}\cup I_{B_4}$ is non-full in $\m B$ together with Claim \ref{abcd}(ii), $I_{B_4}\subset [m'+1,n]$. Thus, $m'+1\in I_{B_4}$ and $I_{B_4}\cap [2,m']=\emptyset$. Let $I'_{B_4}=I_{B_4}\setminus \{m'+1\}$ and 
Let $R_4:=\{1\}\cup I_{B_4}$. 
Define
\begin{equation*}
    \m F_4:=\{F\in \m B(\bar1): I_{B_4}\cap F=\emptyset\}, \,\,\,
    \m G_4:=\m F_{{R_4}}^b\setminus \m B.
\end{equation*}
%Since $\{1,m'+2\}$ is empty in $\m B$ and $I_{B_4}\setminus \{m'+1\}\subset [m'+2,n]$, ${R_4}\cup \{m'+2\}$ is non-full in $\m B$, and hence $\m G_4\ne \emptyset$.
Since $B_4\in \m F_4$, $\m F_4\ne \emptyset$.  
Since $R_4$ is non-full in $\m B$,  $\m G_4\ne \emptyset$.
For any $F\in \m F_4$ and $G\in \m G_4$, there exists a pair $(U,V)$ such that $S_{U,V}(F)=G$. Among all such pairs $(U,V)$, choose one with $V$ inclusion-minimal.

\item[Case 5.] Every set in $\m B(\bar{1})\setminus \m T$ contains $m'+2$, and there exists $B\in \m B(\bar{1})\setminus \m T$ such that $\{1\}\cup I_{B}$ is full in $\m B$ and
$\{1,m'+2\}\cup I_{B}\setminus [2,m']$ is non-full in $\m B$.

Fix $B_5$ that satisfies Case 5. 
Let $I'_{B_5}:=I_{B_5}\setminus [2,m']$ and $R_5:=\{1\}\cup I'_{B_5}$.
Define
\begin{equation*}
    \m F_5:=\{F\in \m B(\bar1): I_{B_5}\cap F=\emptyset\}, \,\,\,
    \m G_5:=\m F_{R_5\cup \{m'+2\}}^b\setminus \m B.
\end{equation*}
Since $B_5\in \m F_5$ and $\{m'+2\}\cup R_5$ is non-full in $\m B$, $\m F_5\ne \emptyset$ and $\m G_5\ne \emptyset$. 
For any $F\in \m F_5$ and $G\in \m G_5$, there exists a pair $(U,V)$ such that $S_{U,V}(F)=G$. Among all such pairs $(U,V)$, choose one with $V$ inclusion-minimal.


\item[Case 6.] Every set in $\m B(\bar{1})\setminus \m T$ contains $m'+2$, and there exists $B\in \m B(\bar{1})\setminus \m T$ such that both $\{1\}\cup I_{B}$ and
$\{1,m'+2\}\cup I_{B}\setminus [2,m']$ are full in $\m B$, and $\{1\}\cup I_{B}\setminus [2,m']$ is non-full in $\m B$.

Fix $B_6$ that satisfies Case 6. 
Let $I'_{B_6}:=I_{B_6}\setminus [2,m']$ and $R_6:=\{1\}\cup I'_{B_6}$.
Define
\begin{equation*}
    \m F_6:=\{F\in \m B(\bar1): I_{B_6}\cap F=\emptyset\}, \,\,\,
    \m G_6:=\m F_{R_6}^b\setminus \m B.
\end{equation*}
Since $B_6\in \m F_6$, $\m F_6\ne \emptyset$. Since $R_6$ is non-full in $\m B$, $\m G_6\ne \emptyset$. 
For any $F\in \m F_6$ and $G\in \m G_6$, there exists a pair $(U,V)$ such that $S_{U,V}(F)=G$. Among all such pairs $(U,V)$, choose one with $V$ inclusion-minimal.

\item[Case 7.] Every set in $\m B(\bar{1})\setminus \m T$ contains $m'+2$, and there exists $B\in \m B(\bar{1})\setminus \m T$ such that $\{1\}\cup I_{B}\setminus [2,m']$ is full in $\m B$ (so both $\{1\}\cup I_{B}$ and 
$\{1,m'+2\}\cup I_{B}\setminus [2,m']$ are full in $\m B$). 

Fix $B_7$ that satisfies Case 7. 
Let $I'_{B_7}:=I_{B_7}\setminus [2,m'+1]$ and $R_7:=\{1\}\cup I'_{B_7}$.
Define
\begin{equation*}
    \m F_7:=\{F\in \m B(\bar1): I_{B_7}\cap F=\emptyset\}, \,\,\,
    \m G_7:=\m F_{R_7\cup \{m'+2\}}^b\setminus \m B.
\end{equation*}
Since $B_7\in \m F_7$, $\m F_7\ne \emptyset$. 
Since $\{1,m'+2\}$ is empty in $\m B$ and $I'_{B_7}\subset [m'+2, n]$, $I'_{B_7}\cup \{1, m'+2\}=R_7\cup \{m'+2\}$ is empty in $\m B$, $\m G_7\ne \emptyset$. 
For any $F\in \m F_7$ and $G\in \m G_7$, there exists a pair $(U,V)$ such that $S_{U,V}(F)=G$. Among all such pairs $(U,V)$, choose one with $V$ inclusion-minimal.
\end{itemize}




We claim that
\begin{itemize}
    \item[(a)] For each $i\in [7]$, $R_i\subset U$; 
    \item[(b)] In Case 3, Case 5, Case 7, $m'+2\not\in V$; 
    \item[(c)] In all cases, $[2,t+\ell]\subset V$;
    \item[(d)] For each $j\in [7]$, $V\cap I_{B_j}=\emptyset.$
\end{itemize}

We see that (a), (b) and (d)  follow directly from 
the definition of $\m F_i, \m G_i$. 
For each $i\in [7]$, since (d) holds and $I_{B_i}$ is a cover of $\m T$, $\m F_i\subset \m B(\bar 1)\setminus \m T$.
By Observation \ref{obser3} and $1\in R_i\subset U$, (c) holds. 

%\andrey{We apply the chosen $S_{U,V}$-shift to the families $\m A,\m B.$}



{\bf Next, we show that $\m M$ and $\m T$ are stable under the chosen $S_{U,V}$-shift. }

Since sets in $\m F_i$ avoid $I_{B_i},$ which is a cover for $\m T,$ we have that $\m T$ is stable. 
We show that $\m M$ is stable. Take $M\in\m M$.
Since $[2,m+1]\subset M$ and $m'<m$, $[2,m'+2]\subset M $.
In Case 1 and Case 2, $m'+2\in U \cap M$, so $\m M$ is stable. In Case 4, since $\{1,m'+2\}$ is empty in $\m B$ (see Claim \ref{abcd}(iv)) and $\{1,m'+2\}\cup I_{B_4}$ is full in $\m B$, we have $I_{B_4}\cap [2,m'+1]\ne \emptyset$.
On the other hand, since $\{1\}\cup I_{B_4}$ is non-full in $\m B$, Claim \ref{abcd}(ii) yields $I_{B_4}\subset [m'+1,n]$. Thus, $m'+1\in I_{B_4}\subset U$. Then $m'+1\in M \cap U$, and $\m M$ is stable.  
In Case 6, since $\{1,m'+2\}$ is empty in $\m B$ and $\{1,m'+2\}\cup I_{B_6}\setminus [2,m']$ is full in $\m B$, 
we have $m'+1\in I_{B_6}\setminus [2,m']=I'_{B_6}\subset R_6\subset U$. Hence $m'+1\in M \cap U$, and $\m M$ is stable.

It remains to treat Case 3, Case 5 and Case 7 together. Let $i\in \{3,5,7\}$ and assume, for contradiction, that $S_{U,V}(M)\not\in \m A$ for some $M\in \m M$. 
  
Since every set in $\m B(\bar{1})\setminus \m T$ contains $m'+2$ and $I_{B_i}$ is a cover of $\m T$, the maximality of $(\m A, \m B)$ implies that $\{1,m'+2\} \cup I_{B_i}$ is full in $\m A$. 
Thus, once we show that
\begin{equation}\label{subsetsuvm}
    \{1,m'+2\}\cup I_{B_i}\subset S_{U,V}(M),
\end{equation}
we obtain $S_{U,V}(M)\in \m A$, a contradiction. 

We now prove (\ref{subsetsuvm}).
Note that $m'+2\in M$, $m'+2\not\in V$ (by (b)) and $I'_{B_i}\subset R_i\subset U$. We have $\{1,m'+2\}\cup I'_{B_i}\subset S_{U,V}(M)$. 
In Case 3, we have $I_{B_3}\subset R_3 \subset U$, and hence $\{1,m'+2\}\cup I_{B_3}\subset S_{U,V}(M)$. 
For $i\in \{5,7\}$, $I_{B_i}\setminus I'_{B_i}\subset [2,m'+1]\subset M$, and since $V\cap I_{B_i}=\emptyset$ (see (d)), it follows that $I_{B_i}\setminus I'_{B_i}\subset S_{U,V}(M)$. 
Thus $\{1,m'+2\}\cup I_{B_i}\subset S_{U,V}(M)$.
%This proves (\ref{subsetsuvm}). 

{\bf Finally,  we show that $S_{U,V}(\m A)$ and $S_{U,V}(\m B\setminus \m T)\cup \m T$
%\andrey{Why are you writing this instead of $S_{U,V}(\m B)$?} 
are cross-intersecting.} 

Note that the above chosen $(U,V)$ may satisfy neither {\bf P} nor {\bf Q}, so the cross-intersecting property does not follow directly from Lemmas \ref{SUV} and \ref{SUV'}, and an additional argument is required. 

%Now, we are going to prove the cross-intersecting property of $(S_{U,V}(\m A), S_{U,V}(\m B\setminus \m T)\cup \m T)$.

Assume for contradiction that there exist $A'\in S_{U,V}(\m A)$ and $B'\in S_{U,V}(\m B\setminus \m T)\cup \m T$ that are disjoint.
Let $A\in \m A, B\in \m B$ be the original sets corresponding to $A', B'$, respectively. To be precise, $A'=S_{U,V}(A)$, $B'=S_{U,V}(B)$ if $B\in \m B\setminus \m T$, and  $B'=B$ if $B\in \m T$.
Put $V':=A\cap B$. 
By the definition of $S_{U,V}$-shift,
we only need to consider the following two cases:
\begin{itemize}
    \item[(I)] $B\ne B'$, $B'\not\in \m B$ and $A'=A$; 
    \item[(II)] $A\ne A'$, $A' \not\in \m A$ and $B=B'$.
\end{itemize}
Since $A'\cap B'=\emptyset$ and $A'=A$ or $B'=B$ in the above two cases,
we conclude that 
%\andrey{It's better to comment on $V'\subsetneqq V$}
\[
A\cap U=\emptyset, \,\,B\cap U=\emptyset, \,\, V'\subsetneqq V,
\]
to see $V'\subsetneqq V$, it's clear that $V'\subset V$, and for case (I), if $V'=V$, then since $A\cap U=\emptyset$, we have $A'=(A\setminus V)\cup U\not\in \m A$, a contradiction; and similarly for case (II).

We claim that for every $i\in [7]$, 
\begin{equation}\label{bcapi}
   B\cap I_{B_i}\ne \emptyset.
\end{equation}
\begin{proof}[Proof of (\ref{bcapi})]
Let $i\in [7]$, and assume for contradiction that $B\cap I_{B_i}= \emptyset$. 
Since $1\in U$ and $B\cap U=\emptyset$, $B\in \m B(\bar1)$.
We claim that $B\in \m F_i$. Indeed, it is clear for $i\ne 2$, and for Case 2, if $B\not\in \m F_2$, then $m'+2\in B$, and hence $m'+2\in B\cap U$, contradicting $B\cap U=\emptyset$.

As $I_{B_i}$ is a cover of $\m T$, $B\in \m B(\bar1)\setminus \m T$, and hence $[2,t+\ell +1]\subset B$. 
Since $[2,m']\subset A$ and $m'\ge t+\ell$, 
$[2,t+\ell]\subset A$, and hence $[2,t+\ell]\subset V'$.
Since $B\cap U=\emptyset$ and $V'\subsetneqq V$, $S_{U',V'}(B)\not\in \m B$ for any $U'\subset U$ with $|U'|=|V'|$. Note that $R_i\subset U$ and $|R_i|\le \ell+2\le |V'|$. 
%\andrey{You mean  $|R_i|\le \ell+2\le |V'|$?} 
So there exists $U'$ satisfying $R_i\subset U'\subset U$ and $|U'|=|V'|$, moreover, we have $S_{U',V'}(B)\in \m G_i$. To see $S_{U',V'}(B)\in \m G_i$, it is trivial for $i\in \{1,2,4\}$, and for $i\in \{3,5,6,7\}$, $B\in \m B(\bar1)\setminus \m T$ gives $m'+2\in B$, so $S_{U',V'}(B)\in \m G_i$. This contradicts the inclusion-minimality of $V$.   
\end{proof}

Let $i\in [7]$.
By (\ref{bcapi}),
Case 1, Case 3 and Case 4 cannot occur, since in these cases $I_{B_i}\subset R_i\subset U$ and $B\cap U=\emptyset$.
So $i\in \{2,5,6,7\}$. 
Since $B\cap I_{B_i}\ne \emptyset$ and $I'_{B_i}\subset R_i\subset U$, we have $B\cap (I_{B_i}\setminus I'_{B_i})\ne \emptyset$.
If there exists $p\in B\cap (I_{B_i}\setminus I'_{B_i})\cap [2,m']$, then since $[2,m']\subset A$, we have $p\in A\cap B=V'$ and so $V'\cap I_{B_i}\ne \emptyset$. This contradicts $V'\cap I_{B_i}\subset V\cap I_{B_i}= \emptyset$ (see (d)). 
Since $I_{B_i}\setminus I'_{B_i} \subset [2,m'+1]$,
\begin{equation}\label{bcapi=m'+1}
 B\cap (I_{B_i}\setminus I'_{B_i} )=\{m'+1\}.   
\end{equation}
Consequently, $I'_{B_i}=I_{B_i}\setminus [2,m'+1]$, and hence 
$\{1,m'+1,m'+2\}\cup I'_{B_2}$ is full in $\m B$. 

We claim that Case 5 and Case 6 cannot occur, since in these cases $I_{B_i}\setminus I'_{B_i} \subset [2,m']$, contradicting (\ref{bcapi=m'+1}).

By (\ref{bcapi=m'+1}) and $V\cap I_{B_i}=\emptyset$, we conclude 
\begin{equation}\label{5-30}
m'+1\not\in A\,\, m'+1\in B, \,\,m'+1\not\in V.    
\end{equation}

Consider Case 7. In this case, $\{1\}\cup (I_{B_7}\setminus [2,m'])$ is full in $\m B$. 
So $A\cap (I_{B_7}\setminus [2,m'])\ne \emptyset$.
Note that $I'_{B_7}=I_{B_7}\setminus [2,m'+1]$.
Since $I'_{B_7}\subset R_7\subset U$ and $A\cap U=\emptyset$, 
$A\cap I'_{B_7}=\emptyset$. Thus $m'+1\in A$, contradicting (\ref{5-30}).

Consider Case 2. In this case, $\{1,m'+2\}\cup I'_{B_2}=R_2\subset U$. 
If (I) occurs, then $B'=(B \setminus V)\cup U$.  We have 
$\{1,m'+1,m'+2\}\cup I'_{B_2}\subset B'$. 
As $\{1,m'+1,m'+2\}\cup I'_{B_2}$ is full in $\m B$, $B'\in \m B$, a contradiction.
If (II) occurs, then $A'=(A \setminus V)\cup U$. 
We have $A\cap (\{1,m'+2\}\cup I'_{B_2})=\emptyset$.  Also, $m'+1\notin A$ by (\ref{5-30}).  
At the same time, $\{1,m'+1,m'+2\}\cup I'_{B_2}$ is full in $\m B$ and $n>a+b$, which implies that $A$ cannot intersect all these sets.
%$m'+1\not\in A$ gives $m'+2\in A$. However, $m'+2\in U$, contradicting $A\cap U=\emptyset$. 
\end{proof}

\end{document}
